\documentclass[11pt]{article}

\usepackage[a4paper,margin=1in]{geometry}
\usepackage{amsmath,amssymb,amsthm,mathtools}
\usepackage{booktabs}
\usepackage{array}
\usepackage{microtype}
\usepackage[colorlinks=true,allcolors=blue]{hyperref}

\newcommand{\Oa}{\mathbb{O}}
\newcommand{\R}{\mathbb{R}}
\newcommand{\C}{\mathbb{C}}
\newcommand{\ii}{\mathrm{i}}
\newcommand{\tr}{\operatorname{tr}}
\newcommand{\rank}{\operatorname{rank}}
\newcommand{\diag}{\operatorname{diag}}
\newcommand{\Span}{\operatorname{span}}
\newcommand{\norm}[1]{\lVert #1\rVert}
\newcommand{\abs}[1]{\lvert #1\rvert}
\newcommand{\ip}[2]{\langle #1,#2\rangle}

\newtheorem{theorem}{Theorem}[section]
\newtheorem{lemma}[theorem]{Lemma}
\newtheorem{proposition}[theorem]{Proposition}
\newtheorem{corollary}[theorem]{Corollary}
\theoremstyle{remark}
\newtheorem{remark}[theorem]{Remark}

\title{Weighted Conditional Negativity and Exact Spectral Obstructions\\
for Octonionic Associator Energy Matrices}
\author{Matthew Connelly}
\date{July 2026}

\begin{document}
\maketitle

\begin{abstract}
Let $u_1,\ldots,u_n$ and $H$ be imaginary octonions and define the pairwise associator-energy matrix
\[
 M_{ab}=\norm{[u_a,u_b,H]}^2,
 \qquad [x,y,z]=(xy)z-x(yz).
\]
After fixing a nonzero vacuum direction, the octonion associator reduces to a complex $SU(3)$ cross product.  On the active support this gives
\[
 M_{ab}=4\bigl(r_a r_b-\abs{G_{ab}}^2\bigr),
\]
where $G$ is a complex Gram matrix and $r_a=G_{aa}$.  We prove that the diagonally normalized matrix is conditionally negative definite in the standard sense and that the raw matrix $M$, by congruence, has at most one positive eigenvalue.  Since $M$ is hollow, every nonzero such matrix has exactly one positive eigenvalue and its ordered singular values obey the universal identity
\[
 s_1=\sum_{j=2}^n s_j.
\]
Consequently, an exact geometric spectrum $(1,x,\ldots,x^{n-1})$ can occur only when $1=x+\cdots+x^{n-1}$.  For $n=3$ the unique positive base is the inverse golden ratio.  We also derive the exact least-squares obstruction to any prescribed three-rung base and show that archived numerical floors coincide with the analytic formula.  The result excludes strongly hierarchical geometric singular spectra for this operator class; it does not exclude Froggatt--Nielsen models generally, other octonionic operators, or nonlinear maps from the spectrum to physical masses.  For complete orthonormal seven-frames we record additional tight-frame identities, while leaving exact hexanacci nonattainability as an open algebraic problem.
\end{abstract}

\section{Introduction}

Octonionic and exceptional-algebraic structures have long been investigated as organizing frameworks for fermionic degrees of freedom and generation structure \cite{GunaydinGursey1973,Baez2002}.  A particularly direct construction assigns to generation directions $u_a\in\operatorname{Im}(\Oa)$ and a vacuum direction $H\in\operatorname{Im}(\Oa)$ the symmetric matrix
\begin{equation}
 M_{ab}(U;H)=\norm{[u_a,u_b,H]}^2.
 \label{eq:def-M}
\end{equation}
Because octonions are alternative, the associator is alternating.  Thus $M$ is entrywise nonnegative and hollow.  It is nevertheless not a Gram matrix in the usual positive-semidefinite sense; throughout this paper we call it the \emph{associator-energy matrix}.

A natural spectral question is whether the ordered singular values of \eqref{eq:def-M} can form a geometric ladder
\begin{equation}
 \frac{(s_1,\ldots,s_n)}{s_1}=(1,x,x^2,\ldots,x^{n-1}),
 \qquad 0<x<1.
 \label{eq:ladder}
\end{equation}
Small choices such as $x\sim\lambda^2$ or $x\sim\lambda^4$, with $\lambda$ the Cabibbo parameter, resemble the hierarchical powers used in Froggatt--Nielsen model building \cite{FroggattNielsen1979,Wolfenstein1983}.  Equation \eqref{eq:ladder}, however, is a specific singular-value ansatz for the operator \eqref{eq:def-M}; it is not a definition of the Froggatt--Nielsen mechanism.

The central observation is that \eqref{eq:def-M} has Lorentzian-type inertia for every finite collection of vectors, with no orthonormality assumption.  The strongest statements and their boundaries are summarized in Table~\ref{tab:status}.

\begin{table}[ht]
\centering
\small
\begin{tabular}{>{\raggedright\arraybackslash}p{0.34\textwidth} >{\raggedright\arraybackslash}p{0.20\textwidth} >{\raggedright\arraybackslash}p{0.36\textwidth}}
\toprule
Claim & Status & Boundary \\
\midrule
$SU(3)$ cross-product reduction & analytic, convention fixed & a different octonion convention changes phases and signs, not the norm identity \\
\addlinespace[2pt]
At most one positive eigenvalue and $s_1=\sum_{j>1}s_j$ & analytic theorem & applies to the operator \eqref{eq:def-M}, not to arbitrary octonionic mass operators \\
\addlinespace[2pt]
Unique admissible geometric base & analytic necessary condition & attainability under extra frame constraints is separate \\
\addlinespace[2pt]
Three-rung residual floor & analytic theorem & the archived optimization uses the squared Euclidean residual defined below \\
\addlinespace[2pt]
Archived frame and matrix checks & verified numerical audit & internal and operator-level consistency, not a replacement for the proof \\
\addlinespace[2pt]
Full seven-frame hexanacci miss & numerical obstruction evidence & no formal emptiness certificate is claimed \\
\bottomrule
\end{tabular}
\caption{Claim-status register.}
\label{tab:status}
\end{table}

\section{Octonionic reduction at fixed vacuum}

\subsection{Conventions}

Let $\Oa$ denote the real octonions in the Cayley--Dickson convention $\Oa=\mathbb{H}\oplus\mathbb{H}\ell$, with multiplication inherited from
\[
 (a,b)(c,d)=\bigl(ac-\overline d\,b,\; da+b\overline c\bigr).
\]
Write $1,e_1,\ldots,e_7$ for the resulting orthonormal basis and identify $\operatorname{Im}(\Oa)$ with $\R^7$.  The norm is multiplicative and the algebra is alternative; in particular the associator is a totally alternating trilinear map \cite{Schafer1966,Baez2002}.

For a nonzero $H$, the transitive action of $G_2=\operatorname{Aut}(\Oa)$ on the unit sphere in $\operatorname{Im}(\Oa)$ allows us to rotate $H/\norm H$ to $e_7$.  Trilinearity then contributes only an overall factor $\norm H^2$ to $M$.  The case $H=0$ is trivial.

In the fixed convention, define
\begin{equation}
 z(u)=\bigl(u_1+\ii u_6,\;u_2-\ii u_5,\;u_3-\ii u_4\bigr)\in\C^3.
 \label{eq:zmap}
\end{equation}
By alternation, an $e_7$ component in either $u$ or $v$ produces an
associator with two equal arguments and therefore does not contribute to
$[u,v,e_7]$.

\begin{lemma}[Fixed-vacuum cross-product identity]
For $u,v\in\operatorname{Im}(\Oa)$, the convention \eqref{eq:zmap} gives
\begin{equation}
 z\bigl([u,v,e_7]\bigr)=2\ii\,\overline{z(u)\times z(v)},
 \qquad
 \norm{[u,v,e_7]}^2=4\norm{z(u)\times z(v)}^2.
 \label{eq:cross}
\end{equation}
\end{lemma}

\begin{proof}
Both sides of the first identity are real-bilinear and alternating in $u,v$.  It is therefore enough to check the basis pairs $e_i,e_j$ using the Cayley--Dickson multiplication displayed above.  The resulting components agree with $2\ii\,\overline{z(e_i)\times z(e_j)}$.  The second identity follows because \eqref{eq:zmap} is an isometry from $e_7^\perp\subset\operatorname{Im}(\Oa)$ to $\C^3$ as real inner-product spaces.
\end{proof}

Let $z_a=z(u_a)$, let $Z\in\C^{3\times n}$ have columns $z_a$, and set
\begin{equation}
 G=Z^\dagger Z,
 \qquad G_{ab}=\ip{z_a}{z_b},
 \qquad r_a=G_{aa}=\norm{z_a}^2.
 \label{eq:Gram}
\end{equation}
The complex Lagrange identity gives
\begin{equation}
 \norm{z_a\times z_b}^2
 =\norm{z_a}^2\norm{z_b}^2-\abs{\ip{z_a}{z_b}}^2,
\end{equation}
so that
\begin{equation}
 M_{ab}=4\bigl(r_a r_b-\abs{G_{ab}}^2\bigr).
 \label{eq:M-Gram}
\end{equation}

\section{Weighted conditional negativity and inertia}

The phrase ``conditionally negative definite'' normally refers to nonpositivity on the hyperplane $\sum_a c_a=0$ \cite{Schoenberg1938}.  The raw matrix $M$ need not have that property when the transverse norms $r_a$ differ.  The correct statement is that a diagonal normalization is conditionally negative definite and that $M$ is congruent to it.  This distinction is essential for standard terminology.

Let $I=\{a:r_a>0\}$ be the active support.  If $r_a=0$, then $z_a=0$ and the corresponding row and column of $M$ vanish.  On $I$, define $D=\diag(r_a)$ and
\begin{equation}
 \widetilde M=D^{-1}MD^{-1}.
 \label{eq:Mtilde}
\end{equation}
Equivalently, with $w_a=z_a/\sqrt{r_a}$,
\begin{equation}
 \widetilde M_{ab}=4\left(1-\abs{\ip{w_a}{w_b}}^2\right).
 \label{eq:normalized-CND}
\end{equation}

\begin{theorem}[Conditional negativity and Lorentzian inertia]
The normalized matrix $\widetilde M$ is conditionally negative semidefinite: for every real vector $d$ satisfying $\sum_{a\in I}d_a=0$,
\begin{equation}
 d^T\widetilde M d\le0.
 \label{eq:CND}
\end{equation}
Consequently, the raw matrix $M$ has at most one positive eigenvalue.
\end{theorem}

\begin{proof}
Let $W$ be the $3\times\abs I$ matrix with columns $w_a$, and define the Hermitian matrix
\[
 A=W\diag(d)W^\dagger.
\]
Since $\norm{w_a}=1$,
\[
 \tr A=\sum_a d_a=0.
\]
Furthermore,
\begin{align}
 \frac14 d^T\widetilde M d
 &=\left(\sum_a d_a\right)^2
   -\sum_{a,b}d_a d_b\abs{\ip{w_a}{w_b}}^2 \\
 &=(\tr A)^2-\tr(A^2)
 =-\tr(A^2)\le0,
\end{align}
where the last inequality follows because $A$ is Hermitian.  This proves \eqref{eq:CND}.

A real symmetric matrix nonpositive on a codimension-one subspace has positive inertia index at most one.  Finally,
\[
 M=D\widetilde M D
\]
on the active support, so Sylvester's law of inertia gives the same positive inertia index for $M$ and $\widetilde M$.  Zero-support rows only add zero eigenvalues.
\end{proof}

The same conclusion can be obtained directly without normalization.  For real $c$, set $A=Z\diag(c)Z^\dagger$.  Then
\begin{equation}
 \frac14 c^TMc=(c^Tr)^2-\tr(A^2),
 \label{eq:weighted-form}
\end{equation}
which is nonpositive on the weighted hyperplane $c^Tr=0$.

\begin{corollary}[Universal singular-value sum rule]
Let $s_1\ge s_2\ge\cdots\ge s_n\ge0$ be the singular values of $M$.  If $M\ne0$, then $M$ has exactly one positive eigenvalue and
\begin{equation}
 \boxed{s_1=\sum_{j=2}^n s_j.}
 \label{eq:sum-rule}
\end{equation}
For $M=0$, the identity is trivial.
\end{corollary}

\begin{proof}
The diagonal of $M$ vanishes, so $\tr M=0$.  A nonzero negative-semidefinite matrix has strictly negative trace; therefore a nonzero $M$ must have a positive eigenvalue.  The theorem shows that this eigenvalue is unique.  Write it as $\lambda_+>0$.  All remaining eigenvalues are nonpositive, and trace zero gives
\[
 \lambda_+=\sum_{\lambda_j<0}\abs{\lambda_j}.
\]
Hence $\lambda_+$ is at least as large as the magnitude of every negative eigenvalue.  It is therefore the largest singular value, and the displayed identity is exactly \eqref{eq:sum-rule}.
\end{proof}

\begin{remark}
For $n=3$, the inertia conclusion also follows from Perron--Frobenius theory
for hollow symmetric nonnegative matrices; related inertia constraints are
studied in \cite{CharlesEtAl2013}.  The normalized conditional-negativity
argument is stronger because it applies to arbitrary finite $n$ and uses the
octonionic image structure rather than a dimension-specific matrix theorem.
\end{remark}

\section{Exact obstruction to geometric ladders}

\begin{corollary}[Unique admissible geometric base]
Suppose $M\ne0$ has the normalized geometric singular spectrum \eqref{eq:ladder}.  Then
\begin{equation}
 1=x+x^2+\cdots+x^{n-1}.
 \label{eq:nbonacci}
\end{equation}
For every $n\ge3$, equation \eqref{eq:nbonacci} has a unique solution in $(0,1)$.  Equivalently, $x$ is the reciprocal of the positive $(n-1)$-bonacci constant.  For $n=3$,
\begin{equation}
 x^2+x-1=0,
 \qquad
 x=\frac{\sqrt5-1}{2}=\varphi^{-1}.
 \label{eq:golden}
\end{equation}
\end{corollary}

\begin{proof}
Substitution of \eqref{eq:ladder} into \eqref{eq:sum-rule} gives \eqref{eq:nbonacci}.  The function $f_n(x)=\sum_{k=1}^{n-1}x^k$ is continuous and strictly increasing on $(0,1)$, with $f_n(0)=0$ and $f_n(1)=n-1>1$.  Hence the root is unique.  Equation \eqref{eq:golden} is the $n=3$ specialization.
\end{proof}

Thus the golden ratio is not selected by an optimization landscape.  It is the unique intersection of the geometric-ladder curve with the exact singular-value hyperplane.

\subsection{Exact realization without an orthonormality constraint}

The preceding corollary is a necessary condition.  In the unrestricted vector domain it is also attainable for $n=3$.

\begin{proposition}[Realization of the full three-rung sum-rule line]
For every $r\in[1/2,1]$, there exist $u_1,u_2,u_3\in\operatorname{Im}(\Oa)$ and $H=e_7$ for which the normalized singular spectrum is
\begin{equation}
 (1,r,1-r).
 \label{eq:line}
\end{equation}
In particular, setting $r=\varphi^{-1}$ gives an exact golden ladder.
\end{proposition}

\begin{proof}
Set $t=\sqrt{(1-r)/2}$ and consider the hollow nonnegative matrix
\begin{equation}
 M_r=
 \begin{pmatrix}
 0&r&t\\
 r&0&t\\
 t&t&0
 \end{pmatrix}.
 \label{eq:Mr}
\end{equation}
The vector $(1,-1,0)^T$ has eigenvalue $-r$.  On the orthogonal span of $(1,1,0)^T$ and $(0,0,1)^T$, the remaining eigenvalues are $1$ and $-(1-r)$.  Hence the singular spectrum is \eqref{eq:line}.

It remains to realize the entries by octonionic associators.  Choose mutually orthogonal complex directions and write
\[
 z_1=A(1,0,0),\qquad z_2=B(0,1,0),\qquad z_3=C(0,0,1).
\]
For any three positive off-diagonal targets $p,q,s$, the equations
\[
 4A^2B^2=p,\qquad4A^2C^2=q,\qquad4B^2C^2=s
\]
have the positive solutions
\[
 A^2=\frac12\sqrt{\frac{pq}{s}},\qquad
 B^2=\frac12\sqrt{\frac{ps}{q}},\qquad
 C^2=\frac12\sqrt{\frac{qs}{p}}.
\]
Taking $(p,q,s)=(r,t,t)$ realizes \eqref{eq:Mr} through
\eqref{eq:cross} for $r<1$.  At $r=1$, take $z_3=0$ and choose
$4A^2B^2=1$.
\end{proof}

This construction does not impose that the $u_a$ form an orthonormal Stiefel frame.  Attainability under that additional constraint is a separate question; the archived computations discussed below provide numerical, not analytic, evidence there.

\subsection{Exact three-rung residual floor}

Let a prescribed base $x\in(0,1)$ define the target $(1,x,x^2)$.  For any nonzero three-vector associator-energy matrix, the normalized spectrum is $(1,r,1-r)$ with $r\in[1/2,1]$.  Define the squared residual
\begin{equation}
 \mathcal E_x(r)=(r-x)^2+(1-r-x^2)^2.
 \label{eq:objective}
\end{equation}

\begin{theorem}[Exact prescribed-base obstruction]
For $0<x<1$,
\begin{equation}
 \min_{1/2\le r\le1}\mathcal E_x(r)
 =\frac{(1-x-x^2)^2}{2},
 \qquad
 r_*(x)=\frac{1+x-x^2}{2}.
 \label{eq:floor}
\end{equation}
The minimum vanishes if and only if $x=\varphi^{-1}$.
\end{theorem}

\begin{proof}
Differentiating \eqref{eq:objective} gives
\[
 \mathcal E_x'(r)=4r-2(1+x-x^2),
\]
so the unique stationary point is $r_*(x)$.  Since $0<x<1$, one has $1/2\le r_*(x)\le5/8$, so it lies in the admissible interval.  Substitution gives \eqref{eq:floor}.  The zero condition is $1-x-x^2=0$.
\end{proof}

By the realization proposition, \eqref{eq:floor} is the exact global minimum over unrestricted vectors.  Under additional orthonormal-frame constraints it remains an exact lower bound.

\section{Numerical artifact audit}

The accompanying artifact contains a $120{,}000$-frame random scan with seed $424242$, optimized candidate frames, and prescribed-base least-squares tests.  The audit script independently reconstructs the octonion product using the Cayley--Dickson convention, recomputes every stored associator-energy matrix from its frame, and checks the spectra and analytic floors.

For the three phenomenologically motivated bases recorded in the artifact, Table~\ref{tab:floors} compares the stored squared objective to \eqref{eq:floor}.

\begin{table}[ht]
\centering
\begin{tabular}{lrr}
\toprule
Base $x$ & Stored objective & Analytic floor \\
\midrule
$\lambda^4=0.0025311454040401$ & $0.4974656674842361$ & $0.4974656674842361$ \\
$\lambda^2=0.05031049$ & $0.4485544838120462$ & $0.4485544838120462$ \\
$\lambda=0.2243$ & $0.2630949706090197$ & $0.2630949706090197$ \\
$\varphi^{-1}$ & $3.08\times10^{-33}$ & $0$ to floating precision \\
\bottomrule
\end{tabular}
\caption{The optimized floors are the analytic projection distances, not unexplained optimizer stalls.  Values are squared Euclidean residuals.}
\label{tab:floors}
\end{table}

The archived best random near-geometric candidate has normalized spectrum
\[
 (1,0.6180408423,0.3819591577)
\]
and geometric residual $2.35\times10^{-10}$.  Among frames with residual below $10^{-6}$, the archived base range is $[0.6175868,0.6184808]$.  These observations are expected from the unique intersection proved above: any spectrum near both the sum-rule line and the geometric curve must lie near $\varphi^{-1}$.

The numerical evidence is corroborative only.  The sum rule and floor formula do not depend on random sampling or optimization.

\section{Complete seven-frames and the remaining open problem}

For a full orthonormal frame $u_1,\ldots,u_7$ of $\operatorname{Im}(\Oa)$ with $H=e_7$, let $x_a=\ip{u_a}{e_7}$ and retain the $3\times7$ complex matrix $Z$ of transverse coordinates.  Completeness gives the tight-frame identities
\begin{equation}
 ZZ^\dagger=2I_3,
 \qquad
 ZZ^T=0.
 \label{eq:tight}
\end{equation}
Consequently the $7\times7$ Hermitian Gram matrix $G=Z^\dagger Z$ satisfies
\begin{equation}
 G^2=2G,
 \qquad
 \rank G=3.
 \label{eq:projection}
\end{equation}
Writing
\begin{equation}
 G=I-xx^T+\ii S,
 \end{equation}
with $S$ real gives
\begin{equation}
 S^T=-S,
 \qquad Sx=0,
 \qquad S^2=xx^T-I.
 \label{eq:S-identities}
\end{equation}
These identities show that a complete frame occupies a much smaller real-algebraic subset than an arbitrary collection of seven vectors.  They are therefore plausible sources of additional attainability obstructions beyond the universal sum rule.

Earlier numerical searches found near-geometric seven-rung spectra displaced from the reciprocal hexanacci base, with the residual concentrated in one weak spectral-moment direction.  Equations \eqref{eq:projection}--\eqref{eq:S-identities} localize the appropriate algebraic arena, but they do not by themselves prove that the exact ladder equations have no solution.  We therefore leave the following problem open.

\begin{quote}
\textbf{Open problem.} Determine whether the polynomial system consisting of \eqref{eq:projection}, the orthonormal-frame constraints, and the exact reciprocal-hexanacci spectral moments is empty.  A proof requires an invariant inequality, elimination argument, interval certificate, or rational sum-of-squares/Positivstellensatz certificate.
\end{quote}

No topological argument or numerical optimizer floor is treated here as an emptiness certificate.

\section{Scope for flavor model building}

The mathematical obstruction is exact but narrow.  If physical masses or Yukawa singular values are identified directly with the singular values of \eqref{eq:def-M}, then the model predicts \eqref{eq:sum-rule}.  A strongly hierarchical normalized target $(1,\epsilon,\delta)$ with $\epsilon+\delta\ll1$ lies a finite distance from the allowed line and cannot be generated by this operator.

There are several ways a broader model could evade the theorem, each of which changes the construction rather than the proof:
\begin{enumerate}
 \item replace squared associator norms by signed or interference-sensitive entries;
 \item enlarge the state space and use a Jordan or Freudenthal observable rather than \eqref{eq:def-M};
 \item apply a nonlinear physical map between the eigen/singular spectrum of $M$ and masses;
 \item introduce additional vacuum or spurion sectors whose combined observable is not hollow and nonnegative.
\end{enumerate}
Accordingly, the result is not a no-go theorem for octonionic flavor or for Froggatt--Nielsen theory.  It closes only the direct geometric-singular-ladder route through the pairwise squared-associator operator \eqref{eq:def-M}.

\section{Conclusion}

The fixed-vacuum octonion associator has a simple complex-geometric form.  After diagonal normalization, its pairwise energy matrix is conditionally negative definite; before normalization it is congruent to that matrix and therefore has at most one positive eigenvalue.  Hollowness upgrades this inertia statement to the exact singular-value identity
\[
 s_1=\sum_{j=2}^n s_j.
\]
The inverse golden ratio in the three-rung problem is thus not an emergent numerical preference.  It is the unique positive base compatible with an exact geometric ladder.  Prescribed Cabibbo-power bases have the exact squared residual floor $(1-x-x^2)^2/2$, matching the archived numerical optimizations.

The positive content is a parameter-free spectral signature and a precise design constraint: any viable extension must alter the operator or the physical spectral map.  The complete-seven-frame identities provide a sharper algebraic arena for future work, but exact hexanacci nonattainability remains unproved and is stated as such.

\section*{Reproducibility statement}

The source package includes the audit script and archived result artifact
\begin{center}
\small
\path{verify_results.py}\\
\path{fn_ladder_rigidity_results.json}.
\end{center}
The script reconstructs the canonical Cayley--Dickson product, recomputes each
stored matrix from its frame, verifies orthonormality and stored spectra, and
checks the analytic floor formula.  It is an independent artifact audit; the
proofs in the paper do not rely on floating-point computation.

\begin{thebibliography}{99}

\bibitem{Baez2002}
J.~C. Baez,
``The Octonions,''
\emph{Bulletin of the American Mathematical Society} \textbf{39} (2002), 145--205.

\bibitem{Schafer1966}
R.~D. Schafer,
\emph{An Introduction to Nonassociative Algebras},
Academic Press, New York, 1966.

\bibitem{GunaydinGursey1973}
M. G\"unaydin and F. G\"ursey,
``Quark structure and octonions,''
\emph{Journal of Mathematical Physics} \textbf{14} (1973), 1651--1667.

\bibitem{Schoenberg1938}
I.~J. Schoenberg,
``Metric spaces and positive definite functions,''
\emph{Transactions of the American Mathematical Society} \textbf{44} (1938), 522--536.

\bibitem{CharlesEtAl2013}
Z.~B. Charles, M. Farber, C.~R. Johnson, and L. Kennedy-Shaffer,
``Nonpositive eigenvalues of hollow, symmetric, nonnegative matrices,''
\emph{SIAM Journal on Matrix Analysis and Applications} \textbf{34} (2013), 1384--1400.

\bibitem{FroggattNielsen1979}
C.~D. Froggatt and H.~B. Nielsen,
``Hierarchy of quark masses, Cabibbo angles and CP violation,''
\emph{Nuclear Physics B} \textbf{147} (1979), 277--298.

\bibitem{Wolfenstein1983}
L. Wolfenstein,
``Parametrization of the Kobayashi--Maskawa matrix,''
\emph{Physical Review Letters} \textbf{51} (1983), 1945--1947.

\end{thebibliography}

\end{document}
